\chapter{Multimedia Representation and Perception}

\textbf{Course:} Multimedia Communications — Marco Cagnazzo

\bigskip\hrule\bigskip

\section{Introduction}

Multimedia traffic dominates modern internet: video streaming alone accounts for \textbf{over 80\% of all internet traffic}.

\textbf{Course scope}: design and implement multimedia systems; understand real-world protocols and standards spanning JPEG, MPEG, H.26x through neural-enhanced coding and semantic delivery.

\textbf{Interdisciplinary pillars}: signal processing + information theory + neural networks.

\bigskip\hrule\bigskip

\section{The Human Visual System}

\subsection{The Eye}

Two types of retinal receptors:

\begin{center}
\begin{tabular}{|l|l|l|l|}
\hline
\textbf{Receptor} & \textbf{Count} & \textbf{Function} & \textbf{Conditions} \\ \hline
\textbf{Cones} & 6–7 million (center) & Color, high resolution & Requires high illumination \\ \hline
\textbf{Rods} & 75–150 million & Light intensity, low resolution & Very sensitive in low light \\ \hline
\end{tabular}
\end{center}
\begin{center}\includegraphics[width=0.600\linewidth]{Pasted_image_20260501163251.png}\end{center}
\subsection{Light Perception}

\begin{itemize}
\item Perceived intensity: \textbf{logarithmic function} of physical intensity
\item Physical dynamic range: $\approx 10^{10}$ (100 dB)
\item HVS cannot operate over the full range simultaneously → adapts via \textbf{global lighting changes} (reduced dynamic range at any given moment)
\item Perceived brightness is not a simple function of luminance — context matters
\end{itemize}
\begin{center}
\begin{minipage}{0.500\linewidth}\centering
\includegraphics[width=\linewidth]{Pasted_image_20260522103557.png}\end{minipage}\hfill\begin{minipage}{0.400\linewidth}\centering\includegraphics[width=\linewidth]{Pasted_image_20260522103612.png}\end{minipage}
\end{center}
\subsection{Contrast Sensitivity Function (CSF)}

\begin{calloutbox}{red!70!black}{Contrast Sensitivity Function (CSF)}
The CSF measures HVS ability to detect luminance differences across \textbf{spatial frequencies}.

\begin{itemize}
\item \textbf{Contrast}: difference in luminance between lightest and darkest parts of a pattern (\%)
\item \textbf{Spatial frequency}: measured in \textbf{cycles per degree} of visual angle
\item \textbf{Sensitivity} = $1 / \text{(minimum detectable contrast)}$
\end{itemize}

CSF \textbf{peaks at mid-range frequencies} (≈ 2–5 cycles/degree) and declines at both low and high frequencies.

\textbf{Implication for compression:} high-frequency errors are less visible → can apply coarser quantization to high frequencies.
\end{calloutbox}
\begin{center}\includegraphics[width=0.806\linewidth]{Pasted_image_20260501163356.png}\end{center}
\subsection{Color Perception}

\begin{itemize}
\item Visible light: \textbf{400–700 nm}
\item Cone distribution by color sensitivity:
\begin{itemize}
\item 65\% sensitive to \textbf{red} (\textasciitilde{}575 nm)
\item 33\% sensitive to \textbf{green} (\textasciitilde{}535 nm)
\item 2\% sensitive to \textbf{blue} (\textasciitilde{}445 nm) — but highly sensitive
\end{itemize}
\end{itemize}
\begin{center}\includegraphics[width=0.600\linewidth]{Pasted_image_20260501163417.png}\end{center}

\begin{itemize}
\item \textbf{Tristimulus}: color sensation corresponds to the 3D response vector of the three cone types
\item Any color = weighted combination of three primary colors
\end{itemize}

\subsection{Color Spaces}

\textbf{RGB} — unit cube representation:
\begin{itemize}
\item $(0,0,0)$: Black; $(1,1,1)$: White; main diagonal: Grayscale
\item Primaries at $(1,0,0)$, $(0,1,0)$, $(0,0,1)$
\item Secondaries: Cyan $(0,1,1)$, Magenta $(1,0,1)$, Yellow $(1,1,0)$
\end{itemize}
\begin{center}\includegraphics[width=0.600\linewidth]{Pasted_image_20260501163800.png}\end{center}
\textbf{HSV} — derived by balancing RGB cube on black corner, rotating so black-white diagonal is vertical, projecting:
\begin{itemize}
\item \textbf{Hue} (H): angle 0°–360° around vertical axis → color family
\item \textbf{Saturation} (S): distance from central axis → color purity
\item \textbf{Value} (V): vertical position 0–1 → brightness
\end{itemize}

\textbf{YCbCr} (≈ digital YUV):
\begin{itemize}
\item \textbf{Y}: luminance (same as HSV Value)
\item \textbf{Cb, Cr}: chrominance (Cartesian representation vs. polar S/H in HSV)
\item $Cb \approx B - Y$, $Cr \approx R - Y$
\end{itemize}

\subsection{RGB to YCbCr Conversion}

\begin{calloutbox}{red!70!black}{ITU-R BT.601 Color Conversion}
\[
\begin{bmatrix} Y \\ Cb \\ Cr \end{bmatrix} = \begin{bmatrix} 0.299 & 0.587 & 0.114 \\ -0.1687 & -0.3313 & 0.5 \\ 0.5 & -0.4187 & -0.0813 \end{bmatrix} \begin{bmatrix} R \\ G \\ B \end{bmatrix} + \begin{bmatrix} 0 \\ 128 \\ 128 \end{bmatrix}
\]

\begin{itemize}
\item \textbf{Y weights} reflect HVS sensitivity: green weighted most (0.587), blue least (0.114)
\item \textbf{Cb, Cr offset} (+128) centers the 8-bit range $[0,255]$
\item \textbf{Invertible}: $\mathbf{p}_{RGB} = M^{-1}(\mathbf{p}_{YCbCr} - \mathbf{o})$ — perfect reconstruction
\end{itemize}
\end{calloutbox}

\subsection{Beyond HVS: Machine-Centric Multimedia}

A growing fraction of multimedia content is consumed by \textbf{algorithms} (surveillance, IoT, industrial monitoring), not humans.

\begin{center}
\begin{tabular}{|l|l|l|}
\hline
\textbf{Aspect} & \textbf{Human Visual System} & \textbf{Machine Models} \\ \hline
Priority & Semantic regions (faces), CSF-limited & High-frequency details, local textures \\ \hline
Goal & Minimize perceived distortion & Maximize downstream task accuracy \\ \hline
Compression strategy & Psychovisually lossless & Discard "semantically irrelevant" content \\ \hline
\end{tabular}
\end{center}

\textbf{Task-Oriented Communications}: if the end-user is a machine, extreme compression ratios are possible by discarding visually important but semantically irrelevant information.

\bigskip\hrule\bigskip

\section{Sound Perception}

\subsection{Pure Tone Perception}

A pure tone $x(t) = a\sin(2\pi f_1 t)$ with power $\sigma^2 = a^2/2$ excites \textbf{multiple nerve fibers} simultaneously (power spreading).

\textbf{Filterbank model}: $M$ filters, each corresponding to one nerve fiber $k$:
\begin{itemize}
\item Frequency response: $H_k(f) = A_k(f)\exp^{j\phi_k(f)}$
\item Response to tone at $f_1$: $y_k(t) = a A_k(f_1)\sin[2\pi f_1 t + \phi_k(f_1)]$
\item \textbf{Spreading function}: $S_E(k) = A_k^2(f_1)$ — power distribution across fibers
\end{itemize}

\subsection{Hearing Threshold}
\begin{center}\includegraphics[width=0.600\linewidth]{Pasted_image_20260501163839.png}\end{center}
\begin{itemize}
\item \textbf{Absolute hearing threshold} $S_a(f)$: minimum power for a tone at frequency $f$ to be audible
\item Minimum of $S_a(f)$ at 1–4 kHz → HVS most sensitive in speech frequency range
\item Compression implication: components below $S_a(f)$ can be discarded
\end{itemize}

\subsection{Critical Band}

\begin{calloutbox}{red!70!black}{Critical Band (CB)}
For $N$ sinusoids \textit{close} to frequency $f_1$, the audibility condition is:
\[
\sum_i \sigma_i^2 > S_a(f_1)
\]
(energy integrates within the critical band)

Sinusoids are "near" if they fall within the \textbf{Critical Band} centered at $f_1$. CB width is a function of $f_1$.
\end{calloutbox}
\begin{center}\includegraphics[width=0.806\linewidth]{Pasted_image_20260501163938.png}\end{center}

\subsection{Masking Curves}

\subsubsection{Frequency Masking}

\textbf{Definition:} $S_m(f_0, \sigma^2, f)$ = minimum power of a tone at $f$ to be audible when a masker at $f_0$ with power $\sigma^2 > S_a(f_0)$ is playing.
\begin{center}\includegraphics[width=0.600\linewidth]{Pasted_image_20260501164014.png}\end{center}

\textbf{Properties of $S_m(f_0, \sigma^2, f)$}:
\begin{itemize}
\item Triangular shape around $f_0$; maximum at $f = f_0$
\item \textbf{Masking index}: $S_m(f, \sigma^2, f) - \sigma^2 < 0$ — masked tone need not be louder than masker
\item Asymmetric: slower decrease toward higher frequencies
\item Slope proportional to CB width
\item Slope (toward higher $f$) decreases as masker power $\sigma^2$ increases
\end{itemize}
\begin{center}\includegraphics[width=0.600\linewidth]{Pasted_image_20260522103757.png}\end{center}
\subsubsection{Time Masking}
\begin{center}\includegraphics[width=0.600\linewidth]{Pasted_image_20260501164647.png}\end{center}
\begin{itemize}
\item \textbf{Pre-masking}: 2–5 ms before onset of masker
\item \textbf{Post-masking}: 100–200 ms after masker ends
\end{itemize}

\subsection{Model Application}

Psychoacoustical model enables:
\begin{itemize}
\item Identify \textbf{non-audible signal components}
\item Allow quantization noise up to the masking threshold
\end{itemize}

\textbf{Limitations}:
\begin{itemize}
\item Model valid only for pure tones or narrowband sounds
\item Only 3-way simultaneous interactions modeled
\item Real signals are far more complex → parameters determined \textbf{experimentally} through large-scale listening tests
\end{itemize}

\bigskip\hrule\bigskip

\section{Image and Video Representation}

\subsection{Gray Level Images}

\begin{itemize}
\item Discrete $N \times M$ grid
\item Raster-order scan: position $(m,n)$ → index $k = (n-1)M + m$
\item $f_{n,m} = f_k$
\end{itemize}
\begin{center}\includegraphics[width=0.600\linewidth]{Pasted_image_20260501164756.png}\end{center}
\subsection{Color Representations}

\textbf{RGB}: three gray-scale components $R, G, B$ — one image per channel.
\begin{center}\includegraphics[width=0.600\linewidth]{Pasted_image_20260501164937.png}\end{center}
\textbf{YUV/YCbCr}: one luminance (Y) + two chrominance (U/Cb, V/Cr) components.
\begin{center}\includegraphics[width=0.600\linewidth]{Pasted_image_20260501164947.png}\end{center}
\subsection{Chroma Subsampling}

Sampling scheme notation: $J : a : b$

\begin{itemize}
\item $J$: reference horizontal size (usually 4)
\item $a$: chroma samples on first line of reference pattern
\item $b$: chroma samples on second line
\end{itemize}
\begin{center}\includegraphics[width=0.600\linewidth]{Pasted_image_20260501165005.png}\end{center}

\begin{center}
\begin{tabular}{|l|l|l|}
\hline
\textbf{Scheme} & \textbf{Horiz. Chroma Res.} & \textbf{Vert. Chroma Res.} \\ \hline
\textbf{4:4:4} & Full & Full \\ \hline
\textbf{4:2:2} & 1/2 & Full \\ \hline
\textbf{4:2:0} & 1/2 & 1/2 \\ \hline
\textbf{4:1:1} & 1/4 & Full \\ \hline
\end{tabular}
\end{center}

\subsection{YCbCr 4:2:0}

Most common format in video compression (JPEG, H.264, HEVC, AV1).

\begin{itemize}
\item \textbf{Y}: full resolution
\item \textbf{Cb, Cr}: half resolution horizontally and vertically
\end{itemize}

\textbf{Compression ratio} vs full RGB:
\[
\frac{4+1+1}{4+4+4} = \frac{6}{12} = 0.5 \quad \Rightarrow \quad 50\% \text{ data reduction}
\]

\textbf{Perceptual basis}: HVS is far more sensitive to luminance (brightness) variations than to chrominance (color) variations → reduced Cb/Cr resolution is largely imperceptible.

\subsection{Quantization}

\begin{itemize}
\item Samples represented on a discrete set: \textbf{uniform quantization} (rounding)
\item $L$ = number of levels; $b = \log_2 L$ bits per sample
\item Standard: $b = 8$ bits per component
\begin{itemize}
\item 256 gray levels (8 bpp)
\item $256^3 \approx 16\text{M}$ colors (24 bpp for RGB)
\end{itemize}
\item High Dynamic Range (HDR): 32–64 bits per channel
\end{itemize}

\subsection{Resolution Standards}

\begin{center}
\begin{tabular}{|l|l|l|}
\hline
\textbf{Format} & \textbf{Resolution} & \textbf{Frame Rate} \\ \hline
SECAM & 384×576 & 50 Hz \\ \hline
PAL & 450×576 & 50 Hz \\ \hline
NTSC & 323×486 & 60 Hz \\ \hline
QCIF & 144×176 & — \\ \hline
CIF & 288×352 & — \\ \hline
4CIF & 576×704 & — \\ \hline
SD/PAL & 720×576 & 50 Hz \\ \hline
HD 720p/i & 1280×720 & 50/100 Hz \\ \hline
HD 1080p/i & 1920×1080 & 50/100 Hz \\ \hline
4K UHD & 3840×2160 & up to 120 Hz \\ \hline
8K UHD & 7680×4320 & up to 120 Hz \\ \hline
\end{tabular}
\end{center}

\subsection{Digital Video Representation}

\[
I : (n, m, T, c) \rightarrow x \in \{0, 1, \ldots, 2^b - 1\}
\]

\begin{itemize}
\item $(n,m)$: spatial position
\item $T$: time index (frame)
\item $c$: color component
\item Sequence of digital images with temporal dependencies, 3 color components (RGB or YUV), possible chroma subsampling
\end{itemize}

\bigskip\hrule\bigskip

\section{Compression Principles}

\subsection{Motivation}

\begin{calloutbox}{green!55!black}{HD DVB Uncompressed Data Rate}
\textbf{System:} 1 luminance component $1920\times1080$ + 2 chrominance components $960\times540$, 8-bit quantization, 50 fps

\[
R = (1920\times1080 + 2\times960\times540)\times8\times50 \approx 1.24 \text{ Gbps}
\]

\textbf{2-hour movie storage:} $\approx 1.12$ TB

\textbf{Conclusion:} compression is essential.
\end{calloutbox}

\subsection{Compression Fundamentals}

Three sources of redundancy:

\begin{enumerate}
\item \textbf{Statistical redundancy}
\begin{itemize}
\item \textit{Spatial}: images are locally homogeneous (neighboring pixels are correlated)
\item \textit{Temporal}: successive video frames are highly correlated
\end{itemize}
\end{enumerate}

\begin{enumerate}
\item \textbf{Psychovisual redundancy}
\begin{itemize}
\item Limits of HVS: CSF, masking, contour sensitivity
\item Goal: \textbf{"Visually lossless"} coding — introduce errors below perception threshold
\end{itemize}
\end{enumerate}

\begin{enumerate}
\item \textbf{Semantic redundancy} \textit{(emerging)}
\begin{itemize}
\item Images contain recognizable structures (objects, concepts)
\item Goal: \textbf{"Generative" coding} — transmit tokens/meaning rather than pixels
\end{itemize}
\end{enumerate}

\textit{Modern algorithms should exploit all three.}

\subsection{Lossless and Lossy Algorithms}

\begin{center}
\begin{tabular}{|l|l|l|}
\hline
\textbf{Property} & \textbf{Lossless} & \textbf{Lossy} \\ \hline
Reconstruction & Perfect & Approximate \\ \hline
Basis & Statistical redundancy only & Quantization + psychovisual redundancy \\ \hline
Compression ratio & Small ($T \leq 3$) & High ($T \approx 5$ to $\gg$) \\ \hline
\end{tabular}
\end{center}

\subsection{Basic Tools for Compression}
\begin{center}\includegraphics[width=0.600\linewidth]{Pasted_image_20260501165037.png}\end{center}
\begin{enumerate}
\item \textbf{Transform}: Linear (DCT, Wavelet) or Non-linear (Neural) — concentrate energy into few coefficients
\item \textbf{Prediction}: Spatial or temporal redundancy reduction
\item \textbf{Quantization}: Rate reduction through approximation (\textbf{the only lossy step})
\item \textbf{Entropy Coding}: Residual redundancy reduction (VLC, Arithmetic coding)
\end{enumerate}

\subsection{Rate and Compression Ratio}

\textbf{Compression ratio}:
\[
T = \frac{B_{\text{in}}}{B_{\text{out}}} = \frac{R_{\text{in}}}{R_{\text{out}}}
\]

\textbf{Coding rate}:
\begin{itemize}
\item Image: $R = B_{\text{out}} / (NM)$ [bpp — bits per pixel]
\item Video/audio: $R = B_{\text{out}} / T$ [bps]
\end{itemize}

\begin{center}
\begin{tabular}{|l|l|}
\hline
\textbf{Application} & \textbf{Typical $T$} \\ \hline
Lossless image coding & $T \leq 3$ \\ \hline
Lossy image coding & $T \approx 5 \to ?$ \\ \hline
Lossy video coding & $T \approx 20 \to ?$ \\ \hline
\end{tabular}
\end{center}

\subsection{Quality and Distortion Metrics}

\subsubsection{Non-Perceptual Objective Criteria (NP-OC)}

\textbf{Error image}: $\mathcal{E}(f, \tilde{f}) = f - \tilde{f}$

\begin{calloutbox}{red!70!black}{Mean Square Error (MSE) and PSNR}
\[
\mathcal{D}(f,\tilde{f}) = \frac{1}{NM}\|\mathcal{E}\|^2 = \frac{1}{NM}\sum_{n,m}\mathcal{E}_{n,m}^2
\]

\[
\text{PSNR}(f,\tilde{f}) = 10\log_{10}\!\left(\frac{255^2}{\mathcal{D}(f,\tilde{f})}\right) \quad [\text{dB}]
\]

\textbf{Pros:} simple, analytically derivable, linked to $\mathcal{L}^2$ norm. \textbf{Cons:} does not model HVS — same MSE can yield very different perceived quality.
\end{calloutbox}

\subsubsection{Perceptual Objective Criteria}

\textbf{Weighted PSNR (WPSNR)}: apply a frequency weighting filter $h$ (modeling CSF):
\[
\text{WPSNR}(f,\tilde{f}) = 10\log_{10}\!\left(\frac{255^2}{\mathcal{D}_W}\right), \qquad \mathcal{D}_W = \frac{1}{NM}\|h * \mathcal{E}\|^2
\]
\begin{center}\includegraphics[width=0.600\linewidth]{Pasted_image_20260501165141.png}\end{center}
\begin{calloutbox}{red!70!black}{Structural Similarity Index (SSIM)}
For image blocks $x$ and $y$ (with $\alpha=\beta=\gamma=1$, $C_3=C_2/2$):
\[
\text{SSIM}(x,y) = \frac{(2\mu_x\mu_y + C_1)(2\sigma_{xy} + C_2)}{(\mu_x^2+\mu_y^2+C_1)(\sigma_x^2+\sigma_y^2+C_2)}
\]

Three components:
\begin{itemize}
\item \textbf{Luminance}: $l(x,y) = \frac{2\mu_x\mu_y + C_1}{\mu_x^2+\mu_y^2+C_1}$
\item \textbf{Contrast}: $c(x,y) = \frac{2\sigma_x\sigma_y + C_2}{\sigma_x^2+\sigma_y^2+C_2}$
\item \textbf{Structure}: $s(x,y) = \frac{\sigma_{xy}+C_3}{\sigma_x\sigma_y+C_3}$
\end{itemize}

Image SSIM = average over all blocks. Range $[0,1]$, 1 = identical.
\end{calloutbox}

\textbf{NN-based metrics (LPIPS)}:
\begin{itemize}
\item PSNR/SSIM fail to capture "realism" in highly compressed or AI-generated content
\item \textbf{LPIPS} (\textit{Learned Perceptual Image Patch Similarity}): uses pre-trained NNs (VGG, ViT) as feature extractors; measures distance in feature space rather than pixel space
\item Key insight: a signal can have low PSNR but high perceptual plausibility
\end{itemize}

\subsubsection{Subjective Criteria (SC)}

\begin{itemize}
\item Ground truth for quality evaluation: human observer experiments
\item Expensive, slow, statistically demanding
\item NP-OC often used as proxy (simplicity, geometric interpretation, analytical optimization) — but correlation with SC is imperfect
\end{itemize}

\subsection{Error Perception}
\begin{center}\includegraphics[width=0.600\linewidth]{Pasted_image_20260501165246.png}\end{center}
\begin{center}\includegraphics[width=0.600\linewidth]{Pasted_image_20260501165300.png}\end{center}
\begin{center}\includegraphics[width=0.600\linewidth]{Pasted_image_20260501165307.png}\end{center}
\begin{center}\includegraphics[width=0.600\linewidth]{Pasted_image_20260501165313.png}\end{center}
\textit{Pepper images with different error types (same MSE=16): distributed noise, concentrated block, contour-based, high-frequency, chroma subsampling}

\textit{All images have the same MSE = 16, yet perceived quality varies enormously. SSIM captures this better than MSE.}

\begin{center}
\begin{tabular}{|l|l|l|}
\hline
\textbf{Error Type} & \textbf{MSE} & \textbf{SSIM} \\ \hline
Distributed white noise ($\sigma=4$) & 16 & 0.906 \\ \hline
Concentrated over $100\times100$ pixels & 16 & 0.972 \\ \hline
Concentrated on contours (Sobel filter) & 16 & 0.987 \\ \hline
Noise on high spatial frequencies & 16 & 0.882 \\ \hline
Chroma subsampling & 21.27 & — \\ \hline
\end{tabular}
\end{center}

\textbf{Key observation}: the same MSE produces very different perceptual quality. Contour errors are least visible (SSIM 0.987); high-frequency noise most visible (SSIM 0.882).

\subsection{Video Perception}

\begin{itemize}
\item HVS has both spatial and \textbf{temporal} sensitivity
\item \textbf{Spatio-temporal CSF}: sensitivity depends on both spatial frequency and temporal frequency
\item Spatial and temporal \textbf{masking} effects apply
\item \textbf{Scene changes} cause sudden sensitivity drops (masking → can tolerate more distortion just after a cut)
\end{itemize}
\begin{center}\includegraphics[width=0.600\linewidth]{Pasted_image_20260501165357.png}\end{center}
\subsection{Complexity, Delay and Robustness}

Three additional performance axes (beyond rate/quality):

\begin{itemize}
\item \textbf{Complexity}: limited by real-time constraints, hardware, economic cost
\item \textbf{Delay}: measured at encoder; related to complexity but mainly driven by \textbf{coding order} (how far ahead encoder must look)
\item \textbf{Robustness}: sensitivity of compressed stream to transmission losses
\end{itemize}

\textbf{Design tensions}:

\begin{center}
\begin{tabular}{|l|l|}
\hline
\textbf{Improve} & \textbf{Conflicts with} \\ \hline
↑ Quality & ↓ Rate \\ \hline
↑ Robustness & ↓ Complexity \\ \hline
— & ↓ Delay \\ \hline
\end{tabular}
\end{center}

Modern system design must balance all five: quality, rate, complexity, delay, robustness.

\textbf{Modern quality assessment approach}: report all three:
\begin{itemize}
\item \textbf{PSNR} (simple, pixel-based)
\item \textbf{SSIM} (structure-based)
\item \textbf{LPIPS} (feature-based, perceptual)
\end{itemize}

\subsection{Codec Comparison}

\begin{center}\includegraphics[width=0.600\linewidth]{Pasted_image_20260501165536.png}\end{center}
\begin{center}\includegraphics[width=0.600\linewidth]{Pasted_image_20260501165547.png}\end{center}
\begin{center}\includegraphics[width=0.600\linewidth]{Pasted_image_20260501165556.png}\end{center}
\begin{center}\includegraphics[width=0.600\linewidth]{Pasted_image_20260501165604.png}\end{center}
\begin{center}\includegraphics[width=0.600\linewidth]{Pasted_image_20260501165612.png}\end{center}
\begin{center}\includegraphics[width=0.600\linewidth]{Pasted_image_20260501165619.png}\end{center}
\begin{center}\includegraphics[width=0.600\linewidth]{Pasted_image_20260501165626.png}\end{center}

\begin{center}
\begin{tabular}{|l|l|l|l|l|}
\hline
\textbf{Codec} & \textbf{Year} & \textbf{Size (Bytes)} & \textbf{Ratio} & \textbf{PSNR (dB)} \\ \hline
Original & — & 6,220,800 & 1.0:1 & — \\ \hline
PNG (lossless) & 1996 & 2,923,960 & 2.1:1 & — \\ \hline
4:2:0 only & — & 3,110,400 & 2.0:1 & 43.05 \\ \hline
JPEG & 1991 & 49,795 & \textbf{124.9:1} & 26.89 \\ \hline
JPEG 2000 & — & 50,943 & 122.1:1 & 29.82 \\ \hline
WebP & 2010 & 42,676 & \textbf{145.8:1} & 30.88 \\ \hline
AVIF & 2019 & 48,745 & 127.6:1 & \textbf{32.45} \\ \hline
\end{tabular}
\end{center}

\textit{All lossy codecs at ≈50 KB target. AVIF achieves highest PSNR at comparable size.}

\bigskip\hrule\bigskip

\section{Summary Table}

\begin{center}
\begin{tabular}{|l|l|l|}
\hline
\textbf{Concept} & \textbf{Key Value / Formula} & \textbf{Note} \\ \hline
Cones & 6–7M (fovea center) & Color, high res, bright light \\ \hline
Rods & 75–150M (periphery) & Intensity, low res, dark \\ \hline
CSF peak & 2–5 cycles/degree & Mid spatial freq most sensitive \\ \hline
Cone distribution & 65\% R / 33\% G / 2\% B &  \\ \hline
RGB→Y weight & $0.299R+0.587G+0.114B$ & ITU-R BT.601 \\ \hline
4:2:0 compression & $6/12 = 0.5$ & 50\% vs full RGB \\ \hline
Pre-masking & 2–5 ms & Temporal masking \\ \hline
Post-masking & 100–200 ms & Temporal masking \\ \hline
MSE & $\frac{1}{NM}\sum\mathcal{E}^2$ & Non-perceptual \\ \hline
PSNR & $10\log_{10}(255^2/\mathcal{D})$ [dB] & Non-perceptual \\ \hline
SSIM & $[l\cdot c\cdot s]$ → $[0,1]$ & Structure-based \\ \hline
LPIPS & NN feature distance & Perceptual, learned \\ \hline
HD uncompressed & ≈1.24 Gbps & 1920×1080, 50 fps, 4:2:0 \\ \hline
Lossless ratio & $T\leq3$ & Stat. redundancy only \\ \hline
Lossy image ratio & $T\approx5\to ?$ & + psychovisual redundancy \\ \hline
Lossy video ratio & $T\approx20\to ?$ & + temporal redundancy \\ \hline
\end{tabular}
\end{center}

